% LivePose — Technical Report
\documentclass[11pt,a4paper]{article}

\usepackage[a4paper, top=2.5cm, bottom=2.5cm, left=2cm, right=2cm]{geometry}
\usepackage[english]{babel}

\usepackage{amsmath,amssymb,amsthm}
\usepackage{algorithm}
\usepackage[noend]{algpseudocode}
\usepackage{booktabs}
\usepackage{array}
\usepackage{graphicx}
\usepackage{xcolor}
\usepackage{enumitem}
\usepackage{microtype}
\usepackage{mathtools}

\usepackage{hyperref}

\hypersetup{
  colorlinks=true,
  linkcolor=blue!70!black,
  urlcolor=blue!70!black,
  citecolor=green!60!black,
}

\newcommand{\R}{\mathbb{R}}
\newcommand{\norm}[1]{\left\lVert#1\right\rVert}

\begin{document}
\begin{titlepage}
\thispagestyle{empty}
\begin{center}
  \vspace*{1.2cm}

  \IfFileExists{logo.png}{%
    \includegraphics[width=0.24\textwidth]{logo.png}\\[1.4cm]
  }{%
    \IfFileExists{docs/logo.png}{%
      \includegraphics[width=0.24\textwidth]{docs/logo.png}\\[1.4cm]
    }{}
  }

  {\Large Indian Institute of Technology Bombay\\[0.35cm]}
  {\large SLP Project Report\\[1.35cm]}

  {\Huge\bfseries LivePose: Real-Time Multi-Camera Markerless 3D Pose Estimation\\[0.45cm]}
  {\Large A Technical Report on Algorithms, Contributions, and Design Evolution\\[1.4cm]}

  {\large \textit{Shreekar S Naik}}
  \\
\textbf{  {\large \textit{22B4518}\\[0.9cm]}}

  {\large Guide: Prof. Priyanshu Agarwal\\[1.1cm]}

  {\large GitHub: \href{https://github.com/ShreekarNaik/SLP-LivePose}{github.com/ShreekarNaik/SLP-LivePose}\\[0.35cm]}
  % {\small \href{https://github.com/ShreekarNaik/SLP-LivePose}{Project repository on GitHub}}\\[1.5cm]

  % \vfill
  % {\large May 2026}
\end{center}
\end{titlepage}
\tableofcontents
\newpage

% ============================================================
\section{Introduction and Motivation}
% ============================================================

LivePose is a real-time multi-camera 3D human pose estimation system built from first
principles, inspired by and benchmarked against the FreeMoCap project~\cite{freemocap}.
FreeMoCap demonstrated that markerless motion capture is feasible with consumer webcams
and open-source pose detectors; however, its architecture is oriented towards
\emph{offline} batch processing of pre-recorded video sessions. LivePose extends this
paradigm in two important directions:

\noindent The project repository is available at \url{https://github.com/ShreekarNaik/SLP-LivePose},
and the README in that repository includes the demo videos and run guide for the system.

\begin{enumerate}
  \item \textbf{Live-first streaming}: The pipeline is designed around a continuous
    frame-streaming loop rather than a post-hoc processing step. Videos, when loaded,
    are treated as frame sources playing at their native FPS through the same path as
    live camera feeds.
  \item \textbf{Robust online estimates}: Rather than accumulating many views of each
    pose and batch-refining, LivePose must produce a plausible 3D skeleton from each
    time step independently, using learned priors carried forward from previous frames.
\end{enumerate}

\noindent The core algorithmic contributions over a naive FreeMoCap-style pipeline are:

\begin{itemize}
  \item Confidence-weighted DLT triangulation with adaptive outlier refinement,
  \item Epipolar consistency voting to reject mis-detected camera views,
  \item Reprojection-guided temporal detection selection,
  \item Adaptive One-Euro filtering instead of fixed-gain Kalman smoothing,
  \item A progressive bone-length model with XPBD-based soft constraint solving.
\end{itemize}

% ============================================================
\section{System Overview}
% ============================================================

\subsection{Pipeline Overview}

LivePose processes synchronized multi-camera video streams through a sequence of
stages that progressively refine the pose estimate. The system first detects 2D
human keypoints independently in each camera view, then selects temporally
consistent detections before applying geometric consistency checks across cameras.
The surviving observations are triangulated into 3D joint positions, followed by
temporal smoothing and biomechanical constraint enforcement to stabilize the final
skeleton estimate.

The resulting 3D skeleton at time $t$ is represented as
$(\mathbf{X}_t, \mathbf{v}_t)$, where $\mathbf{X}_t \in \mathbb{R}^{J \times 3}$
contains the world-space coordinates of the $J=17$ COCO joints and
$\mathbf{v}_t \in \{0,1\}^J$ denotes joint validity.

\subsection{Pose Detection Backend}

The 2D detector is YOLO v8-pose~\cite{yolov8}, selected for its real-time throughput
on CPU (nano variant) and GPU/MPS (small variant). The model is trained on COCO and
outputs $J=17$ keypoints per detected person.

Unlike FreeMoCap (which relies on MediaPipe Holistic), LivePose exposes a pluggable
backend interface and supports MediaPipe and RTMPose as drop-in replacements. All
active backends return a list of detections per camera, so the downstream temporal
selection stage receives the same contract regardless of detector:
\begin{equation}
  \text{process\_batch}(\{I_k\}_{k=1}^{C})
  \;\longrightarrow\;
  \bigl\{\{\hat{p}^{(k)}_n\}_{n=1}^{N_k}\bigr\}_{k=1}^{C}
\end{equation}
where $I_k$ is the frame from camera $k$, $C$ is the number of cameras, $N_k$ is the
number of detected persons in camera $k$, and each detection $\hat{p}^{(k)}_n$ is a
pair $(\mathbf{U}^{(k)}_n, \mathbf{c}^{(k)}_n)$ with
$\mathbf{U}^{(k)}_n \in \R^{J \times 2}$ (pixel coordinates) and
$\mathbf{c}^{(k)}_n \in [0,1]^J$ (per-joint confidence scores).

Detections are sorted by descending mean confidence before being handed to the
selection stage.

% ============================================================
\section{Detection Selection: Temporal Consistency at the 2D Level}
% ============================================================

\subsection{Problem}

A na\"ive approach takes the highest-confidence detection in each camera. This fails
when multiple people are present or when a chair-back, bag, or background silhouette
generates a spurious high-confidence detection. The error propagates immediately into
the 3D estimate because all subsequent stages treat the selected detection as ground
truth.

\subsection{Reprojection-Guided Selection}

LivePose maintains per-frame state:
\begin{align*}
  \hat{\mathbf{X}}_{t-1} &\in \R^{J \times 3} \quad \text{(last triangulated 3D skeleton)}\\
  \hat{\mathbf{u}}^{(k)}_{t-1} &\in \R^2 \quad \text{(2D centroid in camera } k\text{ from previous frame)}
\end{align*}

At time $t$, the reference point for camera $k$ is chosen as follows:
\begin{enumerate}
  \item \textbf{Reprojection prior} (preferred): If $\hat{\mathbf{X}}_{t-1}$ exists,
    project its valid joints into camera $k$ via $\mathbf{P}_k$ and compute the
    centroid of all visible projections:
    \[
      \mu^{(k)}_\text{ref} = \frac{1}{|\mathcal{V}|}\sum_{j \in \mathcal{V}}
      \pi_k\!\left(\hat{\mathbf{X}}_{t-1,j}\right)
    \]
    where $\pi_k(\mathbf{X}) = \mathbf{P}_k [X, Y, Z, 1]^\top / w$ is the projective
    map with projection matrix $\mathbf{P}_k = \mathbf{K}_k[\mathbf{R}_k \mid
    \mathbf{t}_k]$, and $\mathcal{V}$ is the set of valid joint indices.
  \item \textbf{2D centroid prior} (fallback): Use $\hat{\mathbf{u}}^{(k)}_{t-1}$
    from the previous frame's selected detection.
  \item \textbf{Cold start}: Select the detection with highest mean confidence.
\end{enumerate}

Given a reference $\mu^{(k)}_\text{ref}$, each candidate detection $n$ is scored:
\begin{equation}
  s_n = \norm{\hat{c}^{(k)}_n - \mu^{(k)}_\text{ref}}_2
        - 0.1 \cdot \bar{c}^{(k)}_n \cdot d_\text{max}
  \label{eq:selection_score}
\end{equation}
where $\hat{c}^{(k)}_n$ is the centroid of detection $n$ (averaged over joints with
$c_j \geq 0.3$), $\bar{c}^{(k)}_n$ is its mean confidence, and $d_\text{max} =
150\,\text{px}$ is the gate threshold. The second term provides a confidence tiebreaker
while keeping distance dominant.

If $s_n > d_\text{max}$ for all candidates, the camera is marked as occluded (no
detection) for this frame, rather than accepting a potentially wrong person. This hard
gate prevents identity switches from polluting the 3D estimate.

\subsection{Effect on Output}

This stage is the primary defence against \emph{person identity switches}. Without it,
a single camera switching to a background detection can corrupt the triangulated 3D
position for the corresponding half of the skeleton. With it, identity is held stable
across frames as long as the subject moves less than $d_\text{max}$ between frames.

% ============================================================
\section{Epipolar Consistency Filtering}
% ============================================================

\subsection{Motivation}

Even after temporal selection, one camera may have its subject partially occluded and
detect a different person. The epipolar geometry between calibrated cameras provides a
purely geometric test: the projection of a real 3D point onto camera $b$ must lie on
the epipolar line defined by the corresponding point in camera $a$.

\subsection{Fundamental Matrix}

For two fully calibrated cameras $(k, l)$ with intrinsics $\mathbf{K}_k$,
$\mathbf{K}_l$ and extrinsics $(\mathbf{R}_k, \mathbf{t}_k)$,
$(\mathbf{R}_l, \mathbf{t}_l)$, the relative pose from $k$ to $l$ is:
\begin{align}
  \mathbf{R}_\text{rel} &= \mathbf{R}_l \mathbf{R}_k^\top \\
  \mathbf{t}_\text{rel} &= \mathbf{t}_l - \mathbf{R}_l \mathbf{R}_k^\top \mathbf{t}_k
\end{align}

The essential matrix is $\mathbf{E} = [\mathbf{t}_\text{rel}]_\times \mathbf{R}_\text{rel}$
where $[\cdot]_\times$ denotes the skew-symmetric cross-product matrix. The fundamental
matrix is then:
\begin{equation}
  \mathbf{F}_{kl} = \mathbf{K}_l^{-\top} \mathbf{E}\, \mathbf{K}_k^{-1}
  \label{eq:fundamental}
\end{equation}
normalised so that $F_{33} = 1$.

\subsection{Symmetric Epipolar Distance}

For corresponding pixel observations $\mathbf{u}_k \in \R^2$ and $\mathbf{u}_l \in
\R^2$, the symmetric epipolar distance is:
\begin{equation}
  d_\text{epi}(\mathbf{u}_k, \mathbf{u}_l)
  = \frac{1}{2}\left(
      \frac{|{\tilde{\mathbf{u}}_l^\top \mathbf{F} \tilde{\mathbf{u}}_k}|}
           {\norm{(\mathbf{F}\tilde{\mathbf{u}}_k)_{1:2}}}
    + \frac{|{\tilde{\mathbf{u}}_l^\top \mathbf{F} \tilde{\mathbf{u}}_k}|}
           {\norm{(\mathbf{F}^\top\tilde{\mathbf{u}}_l)_{1:2}}}
    \right)
  \label{eq:epi_dist}
\end{equation}
where $\tilde{\mathbf{u}} = [u, v, 1]^\top$ is the homogeneous representation and
the numerator is the algebraic epipolar residual.

\subsection{Camera Consistency Voting}

For each ordered pair $(k, l)$ of cameras, the median symmetric epipolar distance is
computed across all $J = 17$ joints:
\[
  d_{kl} = \operatorname{median}_{j \in \mathcal{V}_{kl}} d_\text{epi}(\mathbf{u}^{(k)}_j, \mathbf{u}^{(l)}_j)
\]
where $\mathcal{V}_{kl}$ is the set of joints with valid (high-confidence) detections
in both cameras. The score of camera $k$ is its median across all pairs involving $k$:
\[
  D_k = \operatorname{median}_{l \neq k} d_{kl}
\]

Camera $k$ is dropped from the frame if $D_k > \bar{D} + \tau_\text{epi}$ where
$\bar{D}$ is the group median and $\tau_\text{epi} = 15\,\text{px}$ by default.
At least two cameras are always retained.

\subsection{Effect on Output}

Epipolar filtering removes cameras that, despite passing the 2D temporal gate, are
geometrically inconsistent with the majority. This is particularly important in
cluttered scenes where the temporal proximity heuristic can be fooled by a second
person who happens to be near where the tracked subject was in the previous frame.
Removing such a camera before triangulation avoids systematic triangulation bias
toward an incorrect detection.

% ============================================================
\section{3D Triangulation: Confidence-Weighted DLT}
% ============================================================

\subsection{Linear Triangulation (DLT)}

Given a 3D point $\mathbf{X} = [X, Y, Z]^\top$ and its projection
$\mathbf{u} = [u, v]^\top$ in camera $k$ with projection matrix
$\mathbf{P}_k \in \R^{3 \times 4}$, the projective constraint is:
\[
  \tilde{\mathbf{u}} \times \mathbf{P}_k \tilde{\mathbf{X}} = \mathbf{0}
\]
Expanding the cross product yields two independent linear equations per camera view:
\begin{align}
  (u\,\mathbf{p}^{(k)3} - \mathbf{p}^{(k)1}) \tilde{\mathbf{X}} &= 0 \label{eq:dlt1}\\
  (v\,\mathbf{p}^{(k)3} - \mathbf{p}^{(k)2}) \tilde{\mathbf{X}} &= 0 \label{eq:dlt2}
\end{align}
where $\mathbf{p}^{(k)i}$ is the $i$-th row of $\mathbf{P}_k$.

\subsection{Confidence Weighting}

For joint $j$ observed in cameras $\mathcal{K}_j = \{k : c^{(k)}_j \geq \theta_c\}$
with $|\mathcal{K}_j| \geq 2$, the weighted system matrix is:
\begin{equation}
  \mathbf{A}_j = \begin{bmatrix}
    c^{(k_1)}_j(u^{(k_1)}_j\,\mathbf{p}^{(k_1)3} - \mathbf{p}^{(k_1)1}) \\
    c^{(k_1)}_j(v^{(k_1)}_j\,\mathbf{p}^{(k_1)3} - \mathbf{p}^{(k_1)2}) \\
    \vdots \\
    c^{(k_m)}_j(u^{(k_m)}_j\,\mathbf{p}^{(k_m)3} - \mathbf{p}^{(k_m)1}) \\
    c^{(k_m)}_j(v^{(k_m)}_j\,\mathbf{p}^{(k_m)3} - \mathbf{p}^{(k_m)2})
  \end{bmatrix} \in \R^{2m \times 4}
  \label{eq:dlt_weighted}
\end{equation}

The homogeneous 3D point $\tilde{\mathbf{X}}_j \in \R^4$ is the right singular vector
of $\mathbf{A}_j$ corresponding to its smallest singular value (solved via SVD). The
Euclidean position is recovered by dehomogenisation: $\mathbf{X}_j = \tilde{X}_{1:3}
/ \tilde{X}_4$.

Weighting by confidence causes high-confidence observations to contribute stronger
row norms in $\mathbf{A}_j$, pulling the solution toward them.

\subsection{Iterative Outlier Refinement}

When $|\mathcal{K}_j| > 2$, a second refinement pass rejects cameras with anomalously
large reprojection errors. The reprojection error for camera $k$ and joint $j$ is:
\[
  \epsilon^{(k)}_j = \norm{\pi_k(\mathbf{X}_j) - \mathbf{u}^{(k)}_j}_2
\]
A camera is dropped if $\epsilon^{(k)}_j > \bar{\epsilon}_j + \sigma_\text{out}
\cdot \hat{\sigma}_j$ where $\bar{\epsilon}_j$ and $\hat{\sigma}_j$ are the empirical
mean and standard deviation of reprojection errors across cameras, and
$\sigma_\text{out} = 2.0$. The system is then re-solved with the surviving cameras.

\subsection{Comparison to FreeMoCap}

FreeMoCap uses an unweighted DLT without refinement. By incorporating per-joint
confidence as weights and performing the iterative outlier removal, LivePose reduces
the influence of uncertain joints (e.g.\ occluded wrists) and corrupted views
on the final 3D estimate.

% ============================================================
\section{Post-Triangulation Validation: Anthropometric Bounds}
% ============================================================

Before the triangulated skeleton enters the filtering and constraint stages, it is
validated against hard anthropometric bounds. These bounds encode the known range of
human bone lengths derived from biomechanics literature, expressed in millimetres for
the COCO-17 joint topology (Table~\ref{tab:bounds}).

\begin{table}[ht]
\centering
\caption{Anthropometric bone length bounds (COCO-17, in mm).}
\label{tab:bounds}
\begin{tabular}{@{}llcc@{}}
\toprule
Joint pair (COCO indices) & Segment & Min (mm) & Max (mm) \\ \midrule
(5, 7), (6, 8) & Upper arm & 200 & 400 \\
(7, 9), (8, 10) & Forearm & 180 & 350 \\
(11, 13), (12, 14) & Femur & 350 & 550 \\
(13, 15), (14, 16) & Tibia & 300 & 500 \\
(5, 6) & Shoulder width & 280 & 500 \\
(11, 12) & Hip width & 200 & 380 \\
(5, 11), (6, 12) & Lateral torso & 400 & 650 \\
\bottomrule
\end{tabular}
\end{table}

For a given frame, the \emph{plausibility score} is:
\begin{equation}
  s_\text{plaus} = \frac{\#\{\text{bones within bounds}\}}{\#\{\text{bones with both endpoints valid}\}}
\end{equation}
If $s_\text{plaus} < 0.5$, the current frame's 3D estimate is discarded and the
previous frame $\hat{\mathbf{X}}_{t-1}$ is used in its place (frame repetition). This
prevents physically implausible skeletons from entering the smoothing and constraint
stages, where they would corrupt the learned bone statistics.

% ============================================================
\section{Temporal Smoothing: Adaptive One-Euro Filter}
% ============================================================

\subsection{Why Not Kalman?}

An earlier version of the pipeline used a Kalman filter for
post-triangulation smoothing. The Kalman filter exhibited two primary issues
in this setting:

\begin{enumerate}
  \item \textbf{Lag during fast motion}: The process noise $Q$ must be set high enough
    to track rapid limb movements; with $Q$ too low, the filter lags by multiple
    frames, distorting fast actions.
  \item \textbf{Noise amplification at rest}: With $Q$ high, the filter provides
    little smoothing when the subject is stationary, amplifying triangulation noise.
\end{enumerate}

The One-Euro filter~\cite{casiez2012} resolves this tradeoff through
speed-adaptive cutoff behaviour: the filter applies stronger smoothing at rest
while remaining responsive during rapid motion.

\subsection{Filter Mathematics}

The One-Euro filter operates as a speed-adaptive first-order low-pass filter. For each
scalar dimension $x$ (the filter is applied independently to all $3J = 51$ scalar
components of the skeleton), with a sample at time $t$ of value $x_t$:

\paragraph{Step 1 — Velocity estimation.} The raw finite-difference velocity is
low-pass filtered with a fixed derivative cutoff $f_d = 1\,\text{Hz}$:
\begin{align}
  \dot{x}_t^\text{raw} &= \frac{x_t - \hat{x}_{t-1}}{\Delta t} \\
  \alpha_d &= \frac{1}{1 + \tau_d / \Delta t}, \quad \tau_d = \frac{1}{2\pi f_d} \\
  \hat{\dot{x}}_t &= \alpha_d \dot{x}_t^\text{raw} + (1 - \alpha_d)\hat{\dot{x}}_{t-1}
\end{align}

\paragraph{Step 2 — Adaptive cutoff.} The position cutoff frequency scales with
estimated speed:
\begin{equation}
  f_c(t) = f_{\min} + \beta\,|\hat{\dot{x}}_t|
  \label{eq:adaptive_cutoff}
\end{equation}
where $f_{\min} = 3.0\,\text{Hz}$ (rest smoothing, configurable) and
$\beta = 0.05$ (speed sensitivity, configurable).

\paragraph{Step 3 — Position filtering.}
\begin{align}
  \alpha_c &= \frac{1}{1 + \tau_c / \Delta t}, \quad \tau_c = \frac{1}{2\pi f_c(t)} \\
  \hat{x}_t &= \alpha_c x_t + (1 - \alpha_c)\hat{x}_{t-1}
\end{align}

\subsection{Effect of Parameters}

\begin{itemize}
  \item \textbf{$f_{\min}$ (rest cutoff)}: Controls how much smoothing is applied to
    near-stationary joints. Higher $f_{\min}$ reduces lag but increases noise at rest.
  \item \textbf{$\beta$ (speed coefficient)}: Controls how quickly the filter opens up
    during motion. Higher $\beta$ tracks fast movements better but introduces more
    noise during transitions.
\end{itemize}

At $f_{\min} = 3\,\text{Hz}$ and $\beta = 0.05$, the filter provides $\approx$2
frames of smoothing at rest and near-zero lag for movements exceeding
$\sim$60\,mm/frame.

\subsection{Filter Lifecycle}

Each joint's filter is reset (re-initialised to the current observation) when the
joint transitions from invalid to valid state, preventing transient artefacts where the
filter attempts to smooth from its last valid position to the new observation.

% ============================================================
\section{Bone Length Constraints}
% ============================================================

\subsection{Motivation}

Triangulation errors and detector noise cause the inferred bone lengths to fluctuate
even when the skeleton is not moving. Since human bone lengths are rigid
(approximately constant), any fluctuation is pure error. Enforcing bone-length
consistency both reduces noise and propagates information from well-observed joints
to poorly-observed ones.

\subsection{Progressive Bone Model}

LivePose uses a subject-specific bone model that learns from observations rather than
relying on population averages. For each bone $b$ (there are $B = 16$ skeletal
connections in COCO-17), the model maintains an online Gaussian estimate
$\mathcal{N}(\mu_b, \sigma_b^2)$ of the true bone length, together with a confidence
$\rho_b \in [0,1]$ that determines how strongly the estimate is enforced.

\subsubsection{Online Update with Adaptive Learning Rate}

Bone length observations enter the model only when:
\begin{enumerate}
  \item Both endpoint joints are valid ($v_a = v_b = 1$).
  \item Both endpoints have reprojection error $< 15\,\text{px}$.
  \item The measured length $\ell_b = \norm{\mathbf{X}_a - \mathbf{X}_b}_2$ is within
    the anthropometric bounds from Table~\ref{tab:bounds}.
\end{enumerate}

Given a clean observation $\ell_b$, the update proceeds as follows. A Gaussian gate
prevents large outliers from perturbing the model:
\begin{equation}
  g_b = \exp\!\left(-\frac{1}{2}\left(\frac{\ell_b - \mu_b}{3\sigma_b}\right)^2\right)
  \label{eq:gate}
\end{equation}
(set to $g_b = 1$ until $n_b \geq 2$ observations have been collected). The effective
learning rate decays with the number of observations:
\begin{equation}
  \alpha_b = \text{clip}\!\left(\frac{\alpha_0 \cdot g_b}{1 + n_b \cdot \lambda},\; \alpha_\text{min},\; 1\right)
  \label{eq:lr}
\end{equation}
with $\alpha_0 = 0.3$, $\lambda = 0.05$ (decay), and $\alpha_\text{min} = 0.01$.
The active progressive model uses an exponentially weighted online update for both
mean and variance:
\begin{align}
  \mu_b &\leftarrow (1 - \alpha_b)\mu_b + \alpha_b \ell_b \\
  \sigma_b^2 &\leftarrow (1 - \alpha_b)\bigl(\sigma_b^2 + \alpha_b(\ell_b - \mu_b')^2\bigr)
  \label{eq:ewvar}
\end{align}
where $\mu_b'$ is the pre-update mean. The confidence $\rho_b$ accumulates as
$\rho_b \leftarrow \min(1.0,\; \rho_b + 0.05)$.

\subsubsection{Compliance Schedule}

The compliance $\gamma_b$ (how loosely the constraint is enforced) follows a discrete
schedule based on observation count $n_b$:
\begin{equation}
  \gamma_b = \begin{cases}
    0.8 & n_b < 10 \quad \text{(initialisation — highly flexible)} \\
    0.4 & 10 \leq n_b < 30 \quad \text{(stabilisation)} \\
    0.1 & n_b \geq 30 \quad \text{(locked — rigid enforcement)}
  \end{cases}
  \label{eq:compliance}
\end{equation}

High compliance early on ensures the model doesn't over-constrain before the subject's
bone lengths are reliably estimated. As confidence increases, constraints tighten.

\subsubsection{Corruption Detection}

Before updating the model, a frame-level corruption check is performed. If more than
50\% of observable bones deviate from their model means by more than $3\sigma_b$, the
entire frame is skipped (neither the model nor the bone estimates are updated). This
prevents a single badly-detected frame from corrupting the running statistics.

\subsection{XPBD Constraint Solving}

\subsubsection{From FABRIK to XPBD}

Initially, FABRIK (Forward And Backward Reaching Inverse Kinematics)~\cite{fabrik} was
considered for the IK step. FABRIK is simple and computationally efficient for open
kinematic chains, but COCO-17 is a \emph{branching} skeleton (multiple chains sharing
a common root), and FABRIK does not natively handle closed-loop or branching
structures without ad-hoc merging logic.

LivePose replaces FABRIK with an XPBD-style (extended Position-Based Dynamics)~\cite{xpbd}
constraint solver, which handles arbitrary constraint topologies uniformly and has a
natural compliance parameter that maps directly onto the learned bone model.

\subsubsection{Algorithm}

The XPBD solver operates in the position domain. Let $\mathbf{p}_j \in \R^3$ denote
the position of joint $j$ (initialised to the smoothed triangulated position for valid
joints). For each bone constraint $(a, b)$ with target length $\ell^*_b$ and effective
weight $w_b$, the following correction is applied at each iteration:

\begin{equation}
  \delta_b = \left(\norm{\mathbf{p}_a - \mathbf{p}_b}_2 - \ell^*_b\right)
             \cdot \frac{\mathbf{p}_a - \mathbf{p}_b}{\norm{\mathbf{p}_a - \mathbf{p}_b}_2}
  \label{eq:xpbd_delta}
\end{equation}

The effective weight $w_b$ is computed from the compliance and stiffness of the bone:
\begin{equation}
  w_b = 1 - \text{clip}\!\left(\frac{\gamma_b}{s_b},\; 0,\; 1\right)
  \label{eq:xpbd_weight}
\end{equation}
where $s_b \in \{0.3, 0.7, 1.0\}$ is the bone stiffness (face/head: 0.3,
limbs: 0.7, torso: 1.0). The position update depends on the validity of each endpoint:
\begin{equation}
  \begin{cases}
    \mathbf{p}_a \leftarrow \mathbf{p}_a - \tfrac{1}{2} w_b \,\delta_b \\
    \mathbf{p}_b \leftarrow \mathbf{p}_b + \tfrac{1}{2} w_b \,\delta_b
  \end{cases}
  \text{ if both } a, b \text{ valid}
  \label{eq:xpbd_both}
\end{equation}
\begin{equation}
  \mathbf{p}_b \leftarrow \mathbf{p}_a
    - \frac{\mathbf{p}_a - \mathbf{p}_b}{\norm{\mathbf{p}_a - \mathbf{p}_b}_2} \cdot \ell^*_b
  \quad\text{if only } a \text{ valid}
  \label{eq:xpbd_chain}
\end{equation}
(symmetric case for only $b$ valid). This is iterated $I = 5$ times per frame.

\subsubsection{Effect on Output}

For fully triangulated joints ($v_a = v_b = 1$), XPBD acts as a soft regulariser:
the positions are nudged by a fraction $w_b$ of the length error, distributing the
correction symmetrically. For joints where one endpoint is invalid (occluded), the
valid endpoint acts as an anchor and the invalid one is placed along the bone direction
at the correct length. This effectively provides \emph{kinematic extrapolation} from
observed joints to occluded ones.

\subsubsection{Influence of the Compliance Parameter}

The compliance schedule from Equation~\eqref{eq:compliance} directly modulates how
much of the bone length error is corrected per iteration. During the learning phase
($\gamma = 0.8$), only $20\%$ of the error is applied per half-joint; in the locked
phase ($\gamma = 0.1$ with torso stiffness $s=1.0$), $90\%$ of the error is corrected.
The gradual tightening prevents large corrections during the observation phase, where
the model mean $\mu_b$ may still be noisy.
% ============================================================
\section{Algorithmic Evolution}
% ============================================================

This section summarises the progression of the major algorithmic decisions during
development.

\subsection{Triangulation}

The earliest implementation used a standard unweighted DLT triangulation approach
similar to existing markerless motion capture systems. Confidence weighting was later
introduced to leverage the per-joint confidence estimates produced by the detector.
A reprojection-error refinement pass was subsequently added to reject inconsistent
camera observations before the final triangulation solve.

\subsection{Temporal Filtering}

\paragraph{Phase 1 — No filtering.}
The initial system directly exposed raw triangulation output, which resulted in
visible high-frequency jitter.

\paragraph{Phase 2 — Kalman filtering.}
A per-joint constant-velocity Kalman filter was introduced to stabilize the
skeleton trajectories. Extended Kalman variants were also explored for direct
multi-camera sensor fusion. Although these approaches improved smoothness, they
required careful tuning of motion and measurement noise parameters and introduced
noticeable lag during rapid movement.

\paragraph{Phase 3 — One-Euro filtering.}
The Kalman-based approach was replaced with the adaptive One-Euro filter, which
provided a more practical balance between responsiveness and noise suppression
under varying motion speeds.

\subsection{Bone Constraints}

\paragraph{Phase 1 — Hard projection.}
The earliest constraint stage projected the skeleton onto a rigid anthropometric
model with fixed population-average bone lengths.

\paragraph{Phase 2 — Online statistical modelling.}
Bone lengths were later estimated online using running statistical estimates,
allowing the model to adapt to the observed subject rather than relying on fixed
human proportions.

\paragraph{Phase 3 — Progressive XPBD enforcement.}
The final system introduced the progressive bone model described in Section~7
together with an XPBD-based constraint solver. FABRIK was initially considered
for inverse kinematics, but XPBD proved better suited to the branching structure
of the COCO-17 skeleton and provided more controllable soft-constraint behaviour.

\subsection{Detection Selection}

The earliest detection stage selected the highest-confidence detection independently
in each camera view. This produced frequent identity switches in cluttered scenes
or multi-person environments. The later reprojection-guided selection stage and
epipolar consistency filtering introduced both temporal and geometric consistency
constraints, substantially improving identity stability across frames.
% ============================================================
\section{Factor Analysis: How Each Component Affects the Output}
% ============================================================

Table~\ref{tab:factor_analysis} summarises how each algorithmic component affects
the quality of the output skeleton.

\begin{table}[ht]
\centering
\caption{Factor analysis: each algorithmic component and its failure mode when absent.}
\label{tab:factor_analysis}
\renewcommand{\arraystretch}{1.3}
\begin{tabular}{@{}p{3.5cm}p{3.5cm}p{3.5cm}p{3cm}@{}}
\toprule
\textbf{Component} & \textbf{Error if absent} & \textbf{How modelled} & \textbf{Key parameter} \\ \midrule
Confidence weighting in DLT & Occluded joints pull 3D position toward occluder & Confidence as row-weight in SVD & $\theta_c = 0.3$ \\
Iterative DLT refinement & Single bad camera biases all joints & $2\sigma$ reprojection outlier rejection & $\sigma_\text{out} = 2.0$ \\
Temporal detection selection & Person identity switches corrupt skeleton & Distance to reprojected prior & $d_\text{max} = 150\,\text{px}$ \\
Epipolar filtering & Wrong-person detection in one camera warps 3D & Symmetric epipolar line distance & $\tau_\text{epi} = 15\,\text{px}$ \\
Anthropometric gate & Physically impossible skeletons enter smoothing & Fraction of bones within bounds & Score $\geq 0.5$ \\
One-Euro filter & High-frequency triangulation noise (rest); lag (motion) & Adaptive cutoff $f_c = f_{\min} + \beta|\dot{x}|$ & $f_{\min} = 3\,\text{Hz}$, $\beta = 0.05$ \\
Bone model learning & Constraints based on wrong lengths distort posture & Exponentially weighted online update with Gaussian gate & $\alpha_0 = 0.3$, $\lambda = 0.05$ \\
XPBD IK & Bone lengths fluctuate; occluded joints not inferred & Soft compliance-weighted position correction & $I = 5$ iterations \\
\bottomrule
\end{tabular}
\end{table}

\subsection{Cascade Effects}

The stages are not independent. Detection selection errors at the 2D level propagate
through the entire chain; a wrong detection in two cameras will produce a triangulation
that is geometrically consistent (it passes the epipolar filter) but anatomically
impossible (caught by the anthropometric gate). Conversely, a clean detection that
produces a high-reprojection-error view (perhaps due to calibration error) is handled
by the DLT refinement pass without the upstream stages needing to intervene.

% ============================================================
\section{Calibration}
% ============================================================

\subsection{Charuco Board}

Camera calibration uses a Charuco board (ChArUco = Chessboard + ArUco), which
combines the robustness of ArUco markers (enabling corner identification even under
partial occlusion) with the sub-pixel accuracy of chessboard corners.

\subsection{Intrinsics}

For each camera $k$, intrinsic calibration solves:
\begin{equation}
  \mathbf{K}_k = \begin{bmatrix} f_x & 0 & c_x \\ 0 & f_y & c_y \\ 0 & 0 & 1 \end{bmatrix},
  \quad \mathbf{d}_k \in \R^5 \quad \text{(distortion coefficients)}
\end{equation}
using the OpenCV Charuco calibration routine, which minimises reprojection error.
The live collection workflow waits for at least 15 board views per camera before
enabling calibration; the lower-level routine itself requires at least 5 valid
Charuco views.

\subsection{Extrinsics}

Extrinsic calibration estimates the relative pose between cameras by identifying
shared board observations. For each synchronized frame containing a board visible to
cameras $k$ and $l$, the board-to-camera rotation and translation
$(\mathbf{R}^{(k)}_\text{board}, \mathbf{t}^{(k)}_\text{board})$ are estimated via
\texttt{cv2.solvePnP}. The relative camera pose is:
\[
  \mathbf{R}_{kl} = \mathbf{R}^{(l)}_\text{board} (\mathbf{R}^{(k)}_\text{board})^\top,
  \quad \mathbf{t}_{kl} = \mathbf{t}^{(l)}_\text{board} - \mathbf{R}_{kl}\mathbf{t}^{(k)}_\text{board}
\]

Multiple shared observations are averaged: rotations via hemisphere-aligned normalized
quaternion averaging to avoid antipodal cancellation, and translations via arithmetic
mean. The live collection workflow waits for at least 8 shared frames per pair before
enabling calibration; the compute step requires at least 3 valid shared solvePnP pose
pairs after detection and pose-estimation failures are discarded.

% ============================================================
\section{Complete Pipeline Algorithm}
% ============================================================

\begin{algorithm}[ht]
\caption{LivePose: Per-Frame 3D Skeleton Estimation}
\label{alg:livepose}
\begin{algorithmic}[1]
\Require Cameras $\{C_k\}_{k=1}^{K}$, frames $\{I_k\}_{k=1}^K$,
         prior state $(\hat{\mathbf{X}}_{t-1}, \hat{\mathbf{v}}_{t-1})$
\Ensure 3D skeleton $(\mathbf{X}_t, \mathbf{v}_t)$ and updated state

\State \textbf{// Stage 1: 2D Detection}
\State $\{\{\hat{p}^{(k)}_n\}\}_{k=1}^K \leftarrow \textsc{YoloBatch}(\{I_k\})$

\State \textbf{// Stage 2: Temporal Detection Selection}
\For{$k = 1$ to $K$}
  \State Compute reference $\mu^{(k)}_\text{ref}$ from $\hat{\mathbf{X}}_{t-1}$ or $\hat{\mathbf{u}}^{(k)}_{t-1}$
  \State Score each detection $n$ via Eq.~\eqref{eq:selection_score}
  \State $p^{(k)} \leftarrow \arg\min_n s_n$ if $s^*_n \leq d_\text{max}$ else $\varnothing$
\EndFor

\State \textbf{// Stage 3: Epipolar Consistency}
\State Compute $\mathbf{F}_{kl}$ for all pairs via Eq.~\eqref{eq:fundamental}
\State $\mathcal{K}_\text{valid} \leftarrow \textsc{EpipolarVoting}(\{p^{(k)}\}, \tau_\text{epi}=15\,\text{px})$

\State \textbf{// Stage 4: Triangulation}
\For{$j = 1$ to $J$}
  \State Collect cameras $\mathcal{K}_j = \{k \in \mathcal{K}_\text{valid} : c^{(k)}_j \geq 0.3\}$
  \If{$|\mathcal{K}_j| \geq 2$}
    \State Build $\mathbf{A}_j$ via Eq.~\eqref{eq:dlt_weighted}; solve $\mathbf{X}_j \leftarrow \textsc{SVD}(\mathbf{A}_j)$
    \State \textbf{[Refine]} Reject cameras with $\epsilon^{(k)}_j > \bar{\epsilon}_j + 2\hat{\sigma}_j$; re-solve
    \State $v_j \leftarrow 1$
  \Else
    \State $v_j \leftarrow 0$
  \EndIf
\EndFor

\State \textbf{// Stage 5: Anthropometric Gate}
\State $s_\text{plaus} \leftarrow \textsc{CheckBounds}(\mathbf{X}_t, \mathbf{v}_t)$
\If{$s_\text{plaus} < 0.5$}
  \State $\mathbf{X}_t \leftarrow \hat{\mathbf{X}}_{t-1}$
\EndIf

\State \textbf{// Stage 6: One-Euro Smoothing}
\State $\mathbf{X}_t \leftarrow \textsc{OneEuro}(\mathbf{X}_t, \mathbf{v}_t, \Delta t)$ \Comment{Eq.~\eqref{eq:adaptive_cutoff}}

\State \textbf{// Stage 7: Bone Model Update}
\State $\textsc{BoneModel.Observe}(\mathbf{X}_t, \mathbf{v}_t)$ \Comment{Exponentially weighted update, Eq.~\eqref{eq:ewvar}}

\State \textbf{// Stage 8: XPBD Bone Constraint Solving}
\For{$i = 1$ to $I=5$}
  \For{each bone $(a, b)$ with an estimated target length}
    \State Compute $\delta_b$ via Eq.~\eqref{eq:xpbd_delta}; apply Eq.~\eqref{eq:xpbd_both} or \eqref{eq:xpbd_chain}
  \EndFor
\EndFor

\State $\hat{\mathbf{X}}_{t} \leftarrow \mathbf{X}_t$; update centroids $\hat{\mathbf{u}}^{(k)}_{t}$
\State \Return $(\mathbf{X}_t, \mathbf{v}_t)$
\end{algorithmic}
\end{algorithm}

% ============================================================
\section{Contributions Summary}
% ============================================================

Table~\ref{tab:contributions} summarises the key contributions over a baseline
FreeMoCap-style offline pipeline.

\begin{table}[ht]
\centering
\caption{Contributions relative to FreeMoCap-style baseline.}
\label{tab:contributions}
\renewcommand{\arraystretch}{1.3}
\begin{tabular}{@{}p{4cm}p{4cm}p{5cm}@{}}
\toprule
\textbf{Component} & \textbf{Baseline (FreeMoCap)} & \textbf{LivePose Contribution} \\ \midrule
Triangulation & Unweighted DLT & Confidence-weighted DLT + reprojection-based outlier refinement \\
Detection & Highest-confidence selection & Reprojection-guided temporal selection with distance gate \\
Cross-camera validation & None & Epipolar consistency voting \\
Temporal filtering & None (batch) & Adaptive One-Euro filter (speed-adaptive cutoff) \\
Bone constraints & None & Progressive online bone model + XPBD soft-constraint IK \\
Validation & None & Anthropometric hard-bounds gate, frame rejection \\
Architecture & Offline batch & Live streaming pipeline with per-frame state \\
\bottomrule
\end{tabular}
\end{table}

% ============================================================
\section{Conclusion}
% ============================================================

LivePose demonstrates that robust online 3D pose estimation from multiple consumer
cameras is achievable by layering complementary consistency constraints at each stage
of the pipeline. The key insight is that no single constraint is sufficient: epipolar
geometry catches wrong detections that temporal proximity misses; anthropometric bounds
catch physically impossible triangulations that the epipolar filter does not flag; and
the XPBD solver enforces bone-length consistency without requiring a fixed human model,
adapting instead to the specific subject observed.

The progression from the initial unweighted DLT through Kalman filtering to the
One-Euro + XPBD combination illustrates a general principle: data-adaptive methods
(speed-adaptive cutoff, online bone learning, compliance scheduling) outperform fixed
parameter methods when the signal statistics are non-stationary, as is inevitable in
real-world motion capture.

% ============================================================
\begin{thebibliography}{9}

\bibitem{freemocap}
Matthis, J.\ S.\ et al.\ (2023).
\textit{FreeMoCap: A free, open-source markerless motion capture system}.
bioRxiv.

\bibitem{yolov8}
Jocher, G.\ et al.\ (2023).
\textit{Ultralytics YOLOv8}.
\url{https://github.com/ultralytics/ultralytics}.

\bibitem{casiez2012}
Casiez, G., Roussel, N., \& Vogel, D.\ (2012).
\textit{1€ Filter: A simple speed-based low-pass filter for noisy input in interactive systems}.
Proceedings of CHI 2012, pp.\ 2527--2530.

\bibitem{fabrik}
Aristidou, A., \& Lasenby, J.\ (2011).
\textit{FABRIK: A fast, iterative solver for the inverse kinematics problem}.
Graphical Models, 73(5), 243--260.

\bibitem{xpbd}
Macklin, M., M\"uller, M., \& Chentanez, N.\ (2016).
\textit{XPBD: Position-based simulation of compliant constrained dynamics}.
Proceedings of MIG 2016, pp.\ 49--54.

\end{thebibliography}

\end{document}
